\documentclass[11pt]{article}
\usepackage[margin=1.1in]{geometry}
\usepackage{amsmath,amssymb,amsthm}
\newtheorem{theorem}{Theorem}\newtheorem{lemma}{Lemma}\newtheorem{conjecture}{Conjecture}
\newtheorem{definition}{Definition}\newtheorem{corollary}{Corollary}\newtheorem*{remark}{Remark}
\title{The Observer--Context Groupoid: Relative Facts,\\ a $\mathbb Z_2$ Switching Invariant, and the Movable Cut}
\author{Manuel Flores Gordillo}
\date{July 2026 --- draft v0.3 (adversarial self-pass; FR theorem added), pending review}
\begin{document}\maketitle
\begin{abstract}
We make ``observer'' a mathematical object: an isometric enlargement of the system
together with a measurement context on the enlarged algebra, organized into a groupoid
carrying the $U(1)$ cocycle of the companion papers. Three results follow, each backed
by machine certificates and one anchored by superconducting-hardware data.
(1)~\emph{Relative facts}: every observer admits an exact local classical section,
while no global value assignment exists for the Peres--Mermin family --- nor for any
enlarged family containing its image under a Wigner enlargement. (2)~A gauge-robust
$\mathbb Z_2$ invariant of switch loops: $+1$ on all loops of Weyl switches, $-1$ on
the oriented mutually-unbiased triangle. Nontriviality \emph{requires} complementary
switching; the converse fails, and the counterexample is pinned. The invariant counts
$\tau$-level switching parity --- the same object as the Wigner--friend $-i$.
(3)~\emph{The movable cut}: the comparison map between cut placements is fixed only up
to section; its ray-level content is invariant (and measured), the section-level
bookkeeping shifts by the loop holonomy, and the update rule available to each observer
is itself forced. (4)~The Frauchiger--Renner chain is \emph{localized}: its observer
cycle has trivial class and exactly invariant holonomy $+1$ --- refuting the natural
conjecture that such paradoxes measure holonomy --- while its contradiction is
state-sector, with contextual fraction exactly $1/6$ and the CHSH inequality at value
$7/3$ as the unique separating facet of the codeword polytope. Observer-dependence, the triangle invariant, and cut conventionality
are thereby derived structures of one cohomological object.
\end{abstract}

\section{Introduction}
The Copenhagen principle of the companion note asserts that quantum theory is locally
valid --- classical within every measurement context --- and globally obstructed. This
paper extends that structure from contexts to \emph{observers}. The blocking issue in
earlier attempts was that ``observer'' remained informal. We repair this with a minimal
definition sufficient for three theorems formalizing the relativity of facts, the
holonomy content of complementary switching, and the movability of the Heisenberg cut.
Every step is pinned by a machine certificate in the public verification suite
(\texttt{observer\_groupoid.py} and companions), and the holonomy split of
Section~\ref{sec:inv} carries a hardware measurement.

\section{The observer--context groupoid}
Fix the finite Weyl model at dimension $d$ with $m$ registers: observables are Weyl
operators $W(v)$, a \emph{context} is a closed commuting Weyl family (not necessarily
maximal), and the section $W(u)W(v)=\tau^{c(u,v)}W(u{+}v)$, $\tau=e^{i\pi/d}$, defines
the $2$-cocycle $c$ of the companion papers.
\begin{definition}[Observer; morphisms]
An \emph{observer} is a pair $o=(\iota,C)$: an isometric enlargement
$\iota:\mathcal H\to\mathcal H\otimes\mathcal H_O$ (the \emph{cut datum}) and a context
$C$ on the enlarged algebra. Morphisms are generated by (i) \emph{switches}: unitaries
of a fixed enlargement carrying $C$ to $C'$, on which the Weyl-section gauge group acts
by left and right multiplication; and (ii) \emph{refinements}: further isometries
$(\iota,C)\to(\iota'\!\circ\iota,\,C\otimes I)$ comparing cut placements. Observers and
these morphisms generate the \emph{observer--context groupoid} $\mathbb O$.
\end{definition}
\begin{lemma}[Pullback and monotonicity]\label{lem:pull}
$A\mapsto A\otimes I_O$ is an injective strict morphism of context structures pulling
the cocycle back on the nose; in particular operator products of embedded contexts are
unchanged, for \emph{any} enlargement. Hence if a family admits no global value
assignment, neither does any enlarged family containing its image: an assignment would
restrict to one for the embedded copy. Certificate: the six embedded Peres--Mermin
context products at $d=8$ are $(+1)^{\times5},(-1)^{\times1}$, identical to $d=4$
(\texttt{observer\_groupoid.py}, part~A).
\end{lemma}

\section{Theorem 1: relative facts}
\begin{theorem}\label{thm:facts}
(a) Every observer admits an exact local classical section: the cocycle restricted to
any single context is a coboundary, with an explicit trivialization $\lambda$ valued in
$\mu_{4d}$; verified at $d=2$ (all six PM contexts), $d=3$, and $d=4$
(\texttt{local\_classicality.py}); half $\tau$-steps are genuinely required.
(b) No global value assignment exists for the PM family, hence none for any enlarged
family containing its image (PM product certificate plus Lemma~\ref{lem:pull}).
\end{theorem}
Each observer's facts are classical; there is no observer-independent assignment of
facts. Observer-dependence is the local-to-global cohomological gap --- a theorem, not
an interpretive gloss. The trivializations $\lambda$ are the observer's classical
bookkeeping, and their mutual incompatibility around obstructed families is the entire
content of contextuality at the observer level.

\section{A $\mathbb Z_2$ switching invariant}\label{sec:inv}
For a loop of switches $U_k\cdots U_1$ closing on its initial context, define the
section holonomy $\mathrm{hol}=\det(W_c\,U_k\cdots U_1)$ with $W_c$ any $SU$ closer.
The value $\mathrm{hol}$ itself is section-dependent (the lesson of the companion
paper's Wigner--friend analysis); its square is not:
\begin{lemma}[Invariance --- proved]\label{lem:inv}
Every Weyl element $w=i^kP$ has $\det(w)^2=1$, and closers are special unitary. Hence
Weyl regauging of any switch multiplies $\mathrm{hol}$ by $\pm1$, and
$\mathrm{hol}^2$ is independent of regauging and of the closer.
(Determinant fact machine-checked over all $16$ Weyl elements.)
\end{lemma}
\begin{theorem}\label{thm:inv}
At $d=2$: (i) every loop of Weyl switches has $\mathrm{hol}^2=+1$;
(ii) the oriented mutually-unbiased triangle $T\to T_x\to T_y\to T$ (switches $H$, $S$)
has $\mathrm{hol}^2=-1$, with $\mathrm{hol}=-i$ in the Clifford presentation. In that
presentation $\mathrm{hol}^2=(-1)^{\#S\bmod 2}$: the invariant is the parity of
$\tau$-level switches. Certificates: $200$ Weyl regaugings and $500$ $SU(2)$ closers,
single values throughout.
\end{theorem}
\begin{corollary}\label{cor:nec}
$\mathrm{hol}^2=-1$ requires complementary switching: loops confined to Weyl switches
(in particular, context-preserving ones) are trivial.
\end{corollary}
\begin{remark}[The converse fails --- pinned]
The $2$-cycle $T\to T_x\to T$ with switches $(H,H)$ traverses a complementary pair yet
has $\mathrm{hol}^2=+1$ (machine-pinned). The invariant therefore detects the
\emph{oriented MUB triangle} --- the triality that produces the Wigner--friend $-i$ ---
not pairwise complementarity. We state this openly; ``complementarity is holonomy'' in
the equivalence sense is false, and necessity (Corollary~\ref{cor:nec}) is the correct
claim.
\end{remark}
Hardware: on \texttt{ibm\_fez} the ray-level holonomy of the triangle loop measures
$+1$ (loop phase $-2.9^\circ\pm1.6^\circ$), excluding $-i$ at ray level by
${\sim}56\sigma$ --- exactly the section split the theorem requires.
\begin{conjecture}
In the Clifford presentation at general $d$, $\mathrm{hol}\in\mu_{2d}$ and the
Weyl-gauge-invariant residue $\mathrm{hol}^2\in\mu_d$ detects the top layer of the
$2$-adic tower of the companion papers.
\end{conjecture}

\section{Theorem 3: the movable cut}
\begin{theorem}\label{thm:cut}
Let two observers differ only by a refinement (friend versus Wigner descriptions of one
run). The comparison map between their descriptions is fixed only up to the section
choice: its ray-level content is invariant --- and is what the hardware measures ---
while the section-level bookkeeping shifts by the loop holonomy of
Section~\ref{sec:inv}. Hence the cut may be moved without empirical consequence, and
every move is cohomologically recorded. Moreover the update rule available to each
observer is not a further convention: statistics, repeatability, and commutant
covariance force exactly the one-parameter L\"uders/depolarizing family per degenerate
block, with the L\"uders instrument the unique efficient member
(\texttt{lueders\_instrument.py}, exact at $(4,[2,2])$ and $(6,[3,2,1])$).
\end{theorem}
This is Copenhagen's conventionality of the cut, derived: the cut is a choice of
trivialization, ray-level physics is invariant under the choice, and the invariant that
survives is precisely the class.

\section{Theorem 4: the Frauchiger--Renner cycle, localized}
Model each lab by one qubit (memory identified with the measured system, as is
standard); the probabilistic core of the FR chain is Hardy's structure on
$\psi=(|h,{\downarrow}\rangle+|t,{\downarrow}\rangle+|t,{\uparrow}\rangle)/\sqrt3$
with the four observer contexts
$\{\bar Z Z\},\{\bar Z X\},\{\bar X X\},\{\bar X Z\}$ --- a closed $4$-cycle
Weyl family in $\mathbb O$ (\texttt{fr\_cycle.py}).
\begin{theorem}\label{thm:fr}
(i) The FR observer cycle has trivial class (all four context phases $+1$; the value
system is solvable, $|L|=16$) and its switch loop closes exactly with
$\mathrm{hol}=+1$; since every two-qubit Weyl element has determinant $+1$, this
holonomy is fully gauge-invariant. The FR paradox is therefore \emph{not} a holonomy
phenomenon: the natural conjecture that observer-cycle paradoxes measure holonomy is
refuted by the machinery of this program itself.
(ii) The paradox is state-sector: $\mathrm{CF}(F_{\mathrm{FR}},\psi)=1/6$ exactly,
and by the codeword-polytope theorem of the companion paper the FR contradiction is
equivalent to the moment vector of $\psi$ lying outside $\mathrm{conv}\,L$. The
separating facet, extracted by exact separation, is the CHSH inequality
$\langle\bar ZZ\rangle-\langle\bar ZX\rangle+\langle\bar XX\rangle
+\langle\bar XZ\rangle\le2$ with quantum value $7/3$, and the
Abramsky--Barbosa--Mansfield normalization reproduces the LP value:
$\mathrm{CF}=(7/3-2)/2=1/6$.
\end{theorem}
The Hardy reading of FR is known; the contribution here is its exact placement in the
program's taxonomy. Consistency conditions for observer networks are polytope-membership
conditions in the state sector, not holonomy conditions: ``absoluteness of observed
events'' fails where the codeword polytope is exited, at a quantified rate.

\section{Discussion and open problems}
Open: the general-$d$ tower refinement of Section~\ref{sec:inv}; multi-round dynamics
beyond a single update; whether any observer-network paradox is genuinely
holonomy-sector (Theorem~\ref{thm:fr} says FR is not --- a candidate would need an
odd cycle, i.e.\ an AvN core); and review of this draft beyond its adversarial
self-passes.
\begin{thebibliography}{9}\small
\bibitem{note} M. Flores Gordillo, \emph{Local validity and contextual holonomy}, note, 2026; companion Papers A--C and verification suite, \texttt{github.com/manuflog/contextuality-obstructions}.
\bibitem{KS} S. Kochen, E. Specker, J. Math. Mech. \textbf{17}, 59 (1967).
\bibitem{Mermin} N. D. Mermin, Rev. Mod. Phys. \textbf{65}, 803 (1993).
\bibitem{AB} S. Abramsky, A. Brandenburger, New J. Phys. \textbf{13}, 113036 (2011).
\bibitem{Wigner} E. P. Wigner, in \emph{The Scientist Speculates} (1961).
\bibitem{FR} D. Frauchiger, R. Renner, Nat. Commun. \textbf{9}, 3711 (2018).
\bibitem{Brukner} \v C. Brukner, Entropy \textbf{20}, 350 (2018).
\bibitem{Bohr} N. Bohr, Phys. Rev. \textbf{48}, 696 (1935).
\bibitem{Heis} W. Heisenberg, \emph{Physics and Philosophy} (1958).
\end{thebibliography}
\end{document}
